% 01-introduction
\section{Introduction}\label{sec:intro}

The goal of the PSIA project is to deploy a realistic, tiered, microservices-based application with AI components on a public cloud, using Infrastructure-as-Code (IaC) end to end: virtual machines are provisioned with Terraform, configured with Ansible, and orchestrated by a \emph{self-hosted} Kubernetes cluster (managed offerings such as GKE are not allowed). The system must also be instrumented and observed with a monitoring stack, and extended with one advanced feature.

Our application is a \emph{Movie Recommendation Hub}. Users register and log in through Keycloak (OpenID Connect~\cite{oidc}), keep a profile with their favourite genres, watched movies and the users they follow, browse and search a catalogue of about 45\,000 movies from The Movies Dataset~\cite{banik2017movies}, and view other users' profiles to discover new interests. An AI inference engine, a small open language model (Qwen2.5~3B) served by Ollama on the VMs' CPUs, together with a sentence-embedding model (\code{nomic-embed-text}) for search by meaning, produces personalised recommendations in two ways: \emph{on demand}, when a user asks for suggestions, optionally with free-text constraints (``a horror movie from the 80s''); and \emph{proactively}, through a background agent that precomputes and caches recommendations per profile, so that users are greeted with fresh suggestions as soon as they log in. Following the database-per-service pattern~\cite{richardson-dbps}, each backend service owns its own PostgreSQL database.

The system runs on Google Cloud Platform (region \code{europe-west4}), on a kubeadm cluster of three VMs (one control-plane node and two workers), and is monitored with Prometheus~\cite{prometheus} and Grafana. As advanced feature we chose \emph{persistent state}: every VM of the deployment can be shut down and restarted without losing any application state.

The rest of this report is organised as follows. Section~\ref{sec:design} presents the architecture, services, data and deployment pipeline. Section~\ref{sec:advanced} describes the persistent-state design and its validation. Section~\ref{sec:eval} studies throughput and latency as the number of concurrent clients grows, and Section~\ref{sec:conclusion} concludes. The demonstration video is available at \url{https://youtu.be/TODO}, and the source code and deployment instructions are in the group repository at \url{https://gitlab.rnl.tecnico.ulisboa.pt/psia/groups26/g17}.
